% ══════════════════════════════════════════════════════════════════════════════
%  EXPERIMENT 5 — THE RABBIT AND FOX MODEL               (Lesson 9, T3 Week 3)
% ══════════════════════════════════════════════════════════════════════════════
\experimentdiv{The Rabbit and Fox Model}{Simulating two populations \textbullet{} A line graph over generations \textbullet{} Judging a model}{9}\label{sec:5}
\lessoncodes{VC2S8U04, VC2S8H03 \textbullet{} VC2S8I04, I05, I06}

\vocabhead
\begin{ctab}
\begin{tabular}{|SW{4.2cm}|W{11.6cm}|}
  \hline\hd & \thead{Word} & \thead{Definition (copy from the slide)} \\ \hline
  & \scaf{\textbf{Population}} & \scaf{All the organisms of one species living in an area at the same time.} \\ \hline
  & \textbf{Predator} & \\ \hline
  & \textbf{Prey} & \\ \hline
  & \textbf{Introduced species} & \\ \hline
  & \textbf{Simulation} & \\ \hline
\end{tabular}
\end{ctab}

\begin{lipanel}
\inquiryli{Run a simulation, graph two populations over time and evaluate the model.}
\understandli{Explain how predator and prey populations depend on each other, using two introduced species.}
\sctitle
\begin{itemize}[leftmargin=1.4em, itemsep=0pt, topsep=1pt]
  \item I can follow the rules of the simulation exactly for 15 generations.
  \item I can draw both populations on one line graph and describe the pattern.
  \item I can say what the model shows well and what it leaves out.
\end{itemize}
\end{lipanel}

\begin{skillsfocus}
\focus{PA1}{Record both populations each generation.}
\focus{PA2}{Draw two lines on one graph, with a key.}
\focus{PA4}{Use a physical model of two populations.}
\focus{PA5}{Describe the pattern in the two lines.}
\focus{PA8}{Evaluate the model against the real world.}
\end{skillsfocus}

% ─── TECHNIQUE ────────────────────────────────────────────────────────────────
\exppart{Technique spotlight}{Blue}

\begin{technique}{Two lines on one graph}{PA2}
When two things are measured at the same times, draw them on the same axes so they can be compared:
\begin{itemize}[leftmargin=1.4em, itemsep=1pt, topsep=3pt]
  \item Time (generation) goes on the horizontal axis.
  \item Use one scale on the vertical axis that fits the larger population.
  \item Use a different mark for each: $\times$ joined by a solid line for rabbits, $\circ$ joined by a dashed line for foxes.
  \item Add a \textbf{key} that says which line is which.
\end{itemize}
Here each point is an exact count, not a measurement with error, so joining the points is the right choice.
\end{technique}

\begin{technique}{Judging a model}{PA8}
A model is simpler than the real thing on purpose. To judge it, ask:
\begin{enumerate}[leftmargin=1.6em, itemsep=1pt, topsep=3pt]
  \item What does it show \textbf{well}?
  \item What does it \textbf{leave out}?
  \item How could it be made \textbf{better}?
\end{enumerate}
\end{technique}

\begin{wedo}{Judging a model}
A globe is a model of Earth. Give one thing it shows well and one thing it leaves out.
\answerbox{3cm}
\end{wedo}

% ─── PLAN AND METHOD ─────────────────────────────────────────────────────────
\exppart{The simulation}{Teal}

\begin{readingbox}{Background --- rabbits and foxes in Australia}
European rabbits were released in Victoria in 1859 and spread across most of the continent within 50 years. Red foxes were released in the 1870s, partly to control them. Both are \textbf{introduced species}: foxes now prey on native animals such as bandicoots as well as rabbits. In this simulation, paper squares stand for rabbits and a card stands for each fox.
\end{readingbox}

\textbf{Aim:} To model how the numbers of rabbits (prey) and foxes (predators) change over 15 generations.

\subsubsection*{Equipment}
An A3 sheet as the ``field'' \textbullet{} 200 paper squares (2\,cm) as rabbits \textbullet{} 20 card squares (6\,cm) as foxes

\subsubsection*{The rules}
\begin{enumerate}[itemsep=2pt]
  \item \textbf{Start:} spread 5 rabbits on the field. Have 1 fox ready.
  \item \textbf{Hunt:} drop each fox onto the field from 30\,cm above it. Every rabbit the fox card touches is eaten: take it off.
  \item \textbf{Foxes:} a fox that caught 3 or more rabbits survives and has one cub (add one fox for next time). A fox that caught fewer than 3 dies.
  \item \textbf{Rabbits:} every rabbit left on the field has one baby: add one rabbit for each. Spread them out. (The field holds at most 200.)
  \item \textbf{Record} the number of rabbits and foxes in Table 5.1. That is one generation.
  \item If all the foxes die, 1 new fox moves in next generation. If all the rabbits are eaten, 3 new rabbits move in.
\end{enumerate}

% ─── RESULTS ─────────────────────────────────────────────────────────────────
\exppart{Results}{Amber}

\subsubsection*{Table 5.1: Rabbits and foxes over 15 generations \skr{PA1}}
\begin{ctab}
\begin{tabular}{|SW{2.6cm}|C{2.2cm}|C{2.2cm}||C{2.6cm}|C{2.2cm}|C{2.2cm}|}
  \hline\hd & \thead{Generation} & \thead{Rabbits} & \thead{Foxes} & \thead{Generation} & \thead{Rabbits} & \thead{Foxes} \\ \hline
  & \scaf{0} & \scaf{5} & \scaf{1} & 8 & & \\ \hline
  & 1 & & & 9 & & \\ \hline
  & 2 & & & 10 & & \\ \hline
  & 3 & & & 11 & & \\ \hline
  & 4 & & & 12 & & \\ \hline
  & 5 & & & 13 & & \\ \hline
  & 6 & & & 14 & & \\ \hline
  & 7 & & & 15 & & \\ \hline
\end{tabular}
\end{ctab}

\needspace{20\baselineskip}
\task Draw both populations on one line graph. Join the points, use a different line for each animal, and add a key and a title. \skr{PA2}
\graphgrid{10cm}

% ─── ANALYSIS ────────────────────────────────────────────────────────────────
\exppart{Analysis and evaluation}{Coral}

\task Describe the pattern of the rabbit line. Use numbers from your table. \skr{PA5}
\answerbox{4cm}

\needspace{10\baselineskip}
\task What happened to the foxes soon after the rabbits increased? Explain why.
\answerbox{4cm}

\needspace{10\baselineskip}
\task What happened to the rabbits after the foxes increased? Explain why.
\answerbox{4cm}

\needspace{14\baselineskip}
\task Judge the model. \skr{PA8}
\begin{ctab}
\begin{tabular}{|W{4.4cm}|W{10.4cm}|}
  \hline\thead{Question} & \thead{Your answer} \\ \hline
  What it shows well & \rule{0pt}{1.6cm} \\ \hline
  What it leaves out (give two) & \rule{0pt}{2.2cm} \\ \hline
  How to make it better & \rule{0pt}{1.6cm} \\ \hline
\end{tabular}
\end{ctab}

\needspace{12\baselineskip}
\task Fox baiting is used in parts of Victoria to protect native animals. Use your model to predict what might happen to the rabbit population if all the foxes in an area were removed, and suggest one problem this could cause for farmers.
\answerbox{5cm}

\begin{yr8box}
\begin{keepext}
\exttask Calicivirus was released in Australia in 1996 and killed many rabbits. Add a rule to the simulation that models a disease, run five more generations with it, and describe what changes for both populations.
\answerbox{6cm}
\end{keepext}
\end{yr8box}

% ─── SKILL CHECK ─────────────────────────────────────────────────────────────
\exppart{Skill check}{NavyBlue}

\begin{skillcheck}
\scrow{PA1}{I recorded both populations every generation.}
\scrow{PA2}{I drew two lines on one graph with a key and a title.}
\scrow{PA4}{I followed the rules of the model exactly.}
\scrow{PA5}{I can describe the pattern in both lines, quoting numbers.}
\scrow{PA8}{I can say what a model shows well, what it leaves out and how to improve it.}
\end{skillcheck}

\task Look back over all five experiments. Which skill have you improved most this unit? Give an example.
\answerbox{5cm}

\end{document}
